\chapter{How to lift a proof without increasing its degree}
\label{tx:chap:lift}
The deletion operation goes from a larger universe to a smaller one.
The lifting operation goes in the other direction: take a low-degree
argument, enlarge the universe by a suitable block, and construct
a larger-universe consequence. These are related viewpoints.
Deletion is an operation on entire invariant spaces and their maps.
The lift is an ordinary linear map on polynomial vectors.

We first specify and verify the lift from the identities of a common
projector. The following technical chapters construct that projector.
Thus the present chapter gives the formula and its mathematical role,
while the later construction proves that an operator with the required
identities exists.

\section{An elementary reason for prime-power steps}
Before the operator construction, consider the simplest counting
phenomenon behind it. Adding an arbitrary number of labels changes
counting expressions. Adding a sufficiently large prime-power block
leaves small binomial counts unchanged in the field.

The basic tool is the \emph{Frobenius map}, $u\mapsto u^p$, in a
commutative ring of characteristic $p$.
Since the intermediate binomial coefficients $\binom pi$ vanish,
$(u+v)^p=u^p+v^p$; multiplication is preserved as well.
Repeating this identity gives the analogous formula for $q=p^a$.
It applies in the polynomial ring $K[T]$ used below, and to the
commuting block-translation operators used later.
Commutation is part of the latter application; no such formula is
being asserted for arbitrary noncommuting matrices.
Integer counts are read in $K$ through the map $m\mapsto m\,1_K$.

\begin{lemma}[Counts below a prime-power block]\label{tx:lem:prime-counts}
Let $K$ have characteristic $p$, let $q=p^a$, and let $n\ge0$.
For every $0\le j<q$, the integer binomial coefficients have the same
image in $K$:
\[
                         \binom{n+q}{j}=\binom nj.
\]
\end{lemma}
\begin{proof}
In $K[T]$, repeated application of the characteristic-$p$ binomial
identity gives $(1+T)^q=1+T^q$. Consequently
\[
                  (1+T)^{n+q}=(1+T)^n(1+T^q).
\]
Multiplication by $T^q$ contributes no coefficient in degrees below
$q$. Comparing the coefficient of $T^j$ proves the assertion.
\end{proof}

For example, in characteristic two, adding four labels preserves
$\binom nj$ for $j=0,1,2,3$. At $n=7$ and $n+4=11$ all four values
are one in the field. Degree four can see the change:
$\binom74=35$ is odd, whereas $\binom{11}4=330$ is even.
This explains why the inequality between block size and complexity
budget is strict.

The count has a polynomial interpretation. For unary variables put
\[
 E_{j,n}=\sum_{1\le i_1<\cdots<i_j\le n}
                           x_{i_1}\cdots x_{i_j}.
\]
Use $E_{0,n}=1$ and $E_{j,n}=0$ for $j>n$.
Its value at the all-one assignment is $\binom nj$.
More strongly, setting the $q$ added variables to one gives, for $j<q$,
\[
 E_{j,n+q}(x_1,\ldots,x_n,1,\ldots,1)
   =\sum_{h=0}^j\binom qh E_{j-h,n}=E_{j,n}.
\]
The intermediate coefficients $\binom qh$, $0<h<q$, vanish in $K$.

This is a complete elementary calculation, but not a lifting theorem
for arbitrary axioms. Setting new variables to one need not preserve
their equations. Even stable coefficient counts do not yet preserve
the intersections which allow PC to multiply a polynomial after
cancellation. The full lift must preserve that subspace calculation.
The lemma explains the choice of step; the common projector below
provides the extra structure needed for proofs.

\section{Low degree, a large starting universe size, and repeated lifting}
Write $N(d)$ for a sufficient starting universe size for degree-$d$ transfer.
If $n\ge N(d)$, one proof with that budget gives proofs at every
size $n+\ell q$, $\ell\ge0$, with one fixed degree bound $d$.
This is already enough to contradict a separate proof that the
minimum degree tends to infinity along that progression. It is the
argument used in Chapter~\ref{tx:ch:running-order}.

The earlier LICS refutation result~\cite[p.~276]{Riis} summarised
estimates $A_{\NS}(d)\le2^{cd}$ and
$A_{\PC}(d)\le2^{2^{cd}}$, extracted there from earlier work of
Kraj\'i\v cek. A degree bounded by a sufficiently small multiple of
$\log n$ or $\log\log n$, respectively, therefore lies in that
historical transfer range. Inverting the starting-size estimates gives
the earlier logarithmic and double-logarithmic gaps.
These estimates concern the historical first-order refutation setting;
we do not use them as inputs to the general target theorem here.

A small budget is easier to fit inside a given universe, but does not
by itself supply a simple construction of the lift. The first complete
construction in this book is the characteristic-two argument of
Chapter~\ref{tx:ch:quadratic}. Its starting universe size is quadratic in $d$.
The later constructions improve this to a linear starting universe size and
cover every positive characteristic. Each stage answers the same
question: how many retained labels are sufficient to transport the
entire degree-$d$ proof calculation?

\section{What we are allowed to keep fixed}
We now specify the operator which makes those repeated lifts possible.

Recall the three support bounds: one variable uses at most $k$ labels,
one axiom index at most $s$, and one target index at most $t$.
A degree-$d$ monomial therefore uses at most $kd$ labels. The number
$r$ below covers all coordinates of the required diagram, rather than
just one monomial; this is why indexed-target detection has a larger
budget than closure transport alone.

Fix a proof degree $d$ and a support budget $r$. The consequence-space
calculation needs $r\ge kd+s$. To combine it with general indexed target
reflection we will use $r=2kd+s+t$. Fix $q=p^a>r$, retain enough labels
for exactness, and set $N=n+q$. Every later iteration may use the same
$d,r,q$. Only the universe size and the particular witness coefficients
change.

\subsection{A numerical trace of the retained-label test}
Take characteristic two, variable support $k=2$, axiom-index support
$s=1$, input-degree bound $D=1$, input stability threshold $n_0=3$,
target-degree bound $D_G=2$,
and proposed proof degree $d=2$. Compare a fully invariant target with
an indexed request whose target-index support is $t=1$.
The final transfer theorem in Chapter~\ref{tx:chap:assemble-transfer}
gives the following sufficient ranges:
\begin{center}
\begin{tabular}{@{}lccc@{}}
\toprule
Request & Budget $r$ & Least $2$-power $q>r$ & Retained universe size $n$\\
\midrule
Invariant & $kd+s=5$ & $8$ & $n\ge\max\{3,1,4,9\}=9$\\
Indexed & $2kd+s+t=10$ & $16$ & $n\ge\max\{3,1,5,19\}=19$\\
\bottomrule
\end{tabular}
\end{center}
For the invariant request, a degree-two proof at size $9$ lifts to
sizes $17,25,\ldots$. For the indexed request, the corresponding
progression starts at $19$ and runs through $35,51,\ldots$.
These are implications from an existing proof, not assertions that
either example request has such a proof. At each step, $n$ is the
\emph{smaller} universe size and $N=n+q$ the larger one. A starting universe size below
the stated range supplies no conclusion from this test.

The count $2r-1$ is one retained pool sufficient for every tuple
length $b\le r$: the individual requirements $2b-1$ are bounded by
$2r-1$. They are not added together as the witness is repaired.
The block size $q$ is additional to that pool. If $r\ge1$, the least
$p$-power strictly above $r$ satisfies $r<q\le pr$; at a $p$-power
budget the upper equality can occur. Once $d$ is fixed, all three
numbers $r,q,2r-1$ stay fixed through the progression.

The lift should not be thought of as replacing every variable by a new
polynomial and then substituting into every product. Its coefficients
are shared across all relevant spaces through a single group-algebra
expression. This gives linear and diagrammatic compatibility, not a
multiplicative substitution law.

\section{The proof can be reconstructed from the closure}
We spell out reconstruction for PC, whose consequence space is the
iterated closure $C_d$. For NS the same construction transports the
single certificate-map image $I_d$ described in
Chapter~\ref{tx:sec:pc-space}.
Suppose an original PC proof produces $f$. The exact consequence-space
characterisation says $f\in C_d(n)$. If a linear lift sends that vector
into $C_d(N)$, there is a legal PC proof at degree $d$ in the larger
universe. One obtains it by storing derivations for bases of closure
stages, expressing the required lifted vectors in those bases, and
concatenating the permitted inference steps.

This reconstruction is finite, but no efficient algorithm or bound on
proof length is asserted. Degree and length are different resources.
The closure lets us explain why a proof exists without forcing every
larger proof to have the same list of intermediate lines as the original.

\section{A lifting operator for bounded-degree consequences}
Fix a proposed proof degree $d\ge1$. Degree zero uses the separate constant-space calculation from Chapter~\ref{tx:sec:pc-space}. Write $W_n=P_{n,\le d}$,
$H_n=P_{n,\le d-1}$ and $X_n=K[\mathcal V_n]$.
Let $V_j(n)$ be the closure stages from \eqref{tx:eq:closure} and write
$C_d(n)$ for their stable union. A proof ending in a polynomial $f$ says exactly that
$f\in C_d(n)$. For the uniform-consequence theorem we must additionally
compare the prescribed targets at different universe sizes.

Choose a support budget $r\ge kd+s$ and add a block of $q=p^a>r$ labels and write $N=n+q$.
The required properties of the linear operator are:
\begin{equation}\label{tx:eq:newcontract}
 \mathsf L:W_n\longrightarrow W_N,\qquad
 \deg\mathsf L(f)\le\deg f,\qquad \mathsf L(1)=1,\qquad
 \mathsf L(V_j(n))\subseteq V_j(N)\quad\hbox{for every }j.
\end{equation}
The degree assertion for zero is interpreted using its harmless zero
image. This contract carries every degree-$d$ consequence to its specified
projector lift and, since $\mathsf L(1)=1$, carries a refutation
upwards. A legal proof is recovered from the larger closure. Iteration reaches $n+\ell q$ for every $\ell\ge0$.
The maps at different steps need not have the same coefficients.

For refutations, full periodicity needs a little more than the forward
contract. It needs
\begin{equation}\label{tx:eq:newiff}
 1\in C_d(n)\quad\Longleftrightarrow\quad1\in C_d(n+q),
\end{equation}
and it needs this at each lower degree budget too. The exact deletion
construction will supply both directions. A forward lift by itself
would not justify all the arguments below about exact recurring values and representatives of T4 peaks. General targets require the same two-way conclusion for their actual coefficient map, obtained by the image-subquotient detection argument in Chapter~\ref{tx:chap:assemble-transfer}.

Let $\iota:W_n\to W_N$ include each retained monomial, and let
$\rho:W_N\to W_n$ discard all other monomials. We also seek
$\rho\mathsf L=\mathrm{id}$. Plain inclusion has this property, but
usually does not carry smaller axioms to consequences of larger axioms.
The sum-axiom example below will show the failure
explicitly. The remedy is to compose inclusion with a common projector.

A map $\mathsf L$ with $\rho\mathsf L=\mathrm{id}$ is a
\emph{section}, or right inverse, of $\rho$. It chooses one larger
representative for each smaller vector, and is injective because
$\mathsf L(f)=0$ implies $f=\rho\mathsf L(f)=0$.
The lifting problem asks for a section which also respects the
consequence spaces. The term will recur later for sections of
coordinate-deletion maps.

The following verification uses Theorem~\ref{tx:thm:deletion} and the common projector constructed in Chapters~\ref{tx:sec:block-algebra}--\ref{tx:sec:odd}. It describes the actual operator and its compatibility with the PC inference rules.

\section{The lift from a common projector}\label{tx:lift-sec:formula}
We now give the map, leaving the construction of its certificate to the
next sections. Temporarily allow a coefficient extension $E/K$ and use
the same notation for the extended polynomial spaces. Section~\ref{tx:lift-sec:descent}
returns the lifting statement to $K$.

Identify the added block $B$ with $A=\Fp^a$, acting on it by translations
and fixing $[n]$. Let $\sigma_1,\ldots,\sigma_a$ be the coordinate
translations. Choose $c_1,\ldots,c_a\in E$ linearly independent over
$\Fp$ and set
\[
 z=\sum_{i=1}^a c_i(\sigma_i-1)\in E[S_N],\qquad z^p=0.
\]
The group algebra $E[S_N]$ consists of linear combinations of label
permutations; it acts on polynomials by the corresponding linear
combinations of variable renamings. The permutations commute on the
block, and $\sigma_i^p=1$, which proves $z^p=0$.

The finite algebraic target is one element $w\in E[S_N]$ for which
\begin{equation}\label{tx:lift-eq:witness}
 u w u=u\quad\text{on every }P_b(N)=E[\Inj([b],[N])],\quad b\le r,
 \qquad u=z^{p-1}.
\end{equation}
Thus the same expression must work on every tuple length up to the
support budget. The identity makes $w$ a partial inverse for the
top power $u=z^{p-1}$ on these spaces; it does not invert the nilpotent
operator $z$. Here $P_b(N)$ is a tuple permutation space; the
polynomial ambient space remains $W_N=P_{N,\le d}$.
For $y=uv$ in $\im u$, the witness identity says $uwy=uwuv=uv=y$.
Thus $wy$ is one preimage of $y$ under $u$; the word \emph{partial}
specifies that this inversion is required on $\im u$ only.
The top power sends vectors to the bottom of length-$p$ nilpotent
chains. Its common witness will let the construction remove those
chains while retaining the length-one directions.

For characteristic two put
\[
                         e=(1-zw)(1-wz).
\]
For general $p$, put $m=p-1$ and
\begin{equation}\label{tx:eq:projector}
 S_0=1,\qquad S_{j+1}=(1-z^jwz^{m-j})S_j\ (0\le j<m),
 \qquad e=(1-z^mw)S_m.
\end{equation}
The two formulas agree at $p=2$. For odd characteristic, the witness
identity is used together with a property of these particular tuple
spaces: their nilpotent chains have length one or $p$. The projector
lemma in Chapter~\ref{tx:ch:prime} proves the formula under that
property; nilpotence $z^p=0$ alone is not sufficient. Once these inputs
are established, the formulas are explicit finite algebra expressions.
The lifting operator is
\begin{equation}\label{tx:lift-eq:actual-lift}
 \boxed{\quad\Lift_{n,N,d}(f)=e\,\iota(f).\quad}
\end{equation}
Equivalently, writing $e=\sum_{g\in S_N}a_g g$,
\[
                    \Lift(f)=\sum_{g\in S_N}a_g\,g\bigl(\iota(f)\bigr).
\]
This explains why degree does not increase: inclusion and every
renaming preserve degree. The difficult point is to choose the
\emph{common} coefficients $a_g$ so that all consequence spaces are
preserved. An arbitrary sum of renamings does not have that property.

\subsection{The two facts that make the formula work}
Let $J_r$ be the ideal of group-algebra elements acting as zero on all
$P_b(N)$ for $b\le r$, and let $\mathcal M_r$ consist of the finite
dimensional representations annihilated by $J_r$. This class contains
all the bounded-support permutation spaces needed by PC, and their
invariant subspaces and quotients. It is the class $\mathcal C_r(N)$
of Chapter~\ref{tx:chap:deletion}, now over the auxiliary field $E$.
Define
\[
               \Del(M)=\ker z_M/(\ker z_M\cap\im z_M).
\]
The construction establishes two facts.
\begin{enumerate}
\item On permutation bases, $\Del$ retains exactly the objects supported
on $[n]$. On polynomial spaces its identification is induced by $\rho$:
\[
 \ker z_{W_N}/(\ker z_{W_N}\cap\im z_{W_N})\simeq W_n.
\]
Stable axiom matrices and the actual multiplication maps are transported
to their smaller-universe counterparts. With a permutation space as
first tensor factor, the identification is natural on the other factor.
\item The common $e$ is an idempotent with $ze=ez=0$ on
$\mathcal M_r$, and for $v\in\ker z$,
\begin{equation}\label{tx:lift-eq:onker}
                             ev=v-z^{p-1}wv.
\end{equation}
It identifies $\Del(M)$ naturally, as a vector space, with $eM$.
Consequently $\Del$ preserves kernels, images and exact sequences in
this class.
\end{enumerate}
These facts are proved in
Chapters~\ref{tx:ch:square-zero}--\ref{tx:ch:prime} from the common
witness and the block action. They are stronger than a count of dimensions. They refer to
the actual maps used in the consequence calculation.
The idempotent property means $e^2=e$: applying the common projection
twice has the same effect as applying it once.

\section{Compatibility with the entire PC closure}\label{tx:lift-sec:verify}
Apply the two facts just stated to the recurrence
\eqref{tx:eq:closure}. The usable axioms form the image of a stable
coefficient matrix, so deletion takes $V_0(N)$ to $V_0(n)$. Suppose
it has already taken $V_j(N)$, with its inclusion in $W_N$, to $V_j(n)$.
The intersection $V_j(N)\cap H_N$ is the kernel of
\[
              V_j(N)\oplus H_N\longrightarrow W_N,
              \qquad (f,h)\longmapsto f-h.
\]
Exactness preserves it. The tensor identification and the actual
multiplication map then preserve the multiplication image, and exactness
preserves its sum with $V_j(N)$. Induction gives
\begin{equation}\label{tx:lift-eq:stage-transport}
       \rho\bigl(V_j(N)\cap\ker z_{W_N}\bigr)=V_j(n)
                         \qquad\text{for every }j.
\end{equation}
The tensor domain is a subspace of $X_N\otimes H_N$, whose support bound
is $kd$. No closure of $\mathcal M_r$ under arbitrary tensor products
is assumed.

\begin{proposition}[The lifting formula preserves and reflects consequences]\label{tx:lift-prop:lift}
Assume the two facts in Section~\ref{tx:lift-sec:formula}. Then
$\Lift=e\iota$ preserves degree, $\rho\Lift=\mathrm{id}$,
$\Lift(1)=1$, and
\[
                  \Lift(f)\in V_j(N)\quad\Longleftrightarrow\quad
                  f\in V_j(n).
\]
The same equivalence holds for the stable spaces $C_d$.
\end{proposition}
\begin{proof}
Every retained monomial is fixed by the block action, so $z\iota(f)=0$.
Equation~\eqref{tx:lift-eq:onker} gives
\[
                  \Lift(f)=\iota(f)-z^{p-1}w\iota(f).
\]
The image of $z$ has no retained component: $z$ kills retained basis
vectors and preserves the span of moved ones. Hence $\rho\Lift(f)=f$.
On constants, every permutation acts trivially, so the correction term
vanishes and $\Lift(1)=1$.

If $f\in V_j(n)$, equation~\eqref{tx:lift-eq:stage-transport} supplies
$v\in V_j(N)\cap\ker z$ with $\rho v=f$. The permutation-space
identification says that $v-\iota(f)\in\ker z\cap\im z$. Since
$e\im z=0$, we have $\Lift(f)=ev$. The space $V_j(N)$ is invariant
under all label permutations, hence under $e$, and $ev\in V_j(N)$.
Conversely, $\Lift(f)\in V_j(N)$ is killed by $z$, so
\eqref{tx:lift-eq:stage-transport} and $\rho\Lift(f)=f$ give $f\in V_j(n)$.
Degree preservation follows from the renaming formula. Pass to the
unions for the stable statement.
\end{proof}

This proves both the forward lifting implication and
\eqref{tx:eq:newiff}. The image of the lift need not be all of $C_d(N)$:
the larger universe has many additional consequences. Every smaller consequence has this projector lift. For $n\ge n_*$ and a fully invariant target whose stabilised degree
is at most $d$, naturality of its map from $K$ gives
$\Lift(G^{(n)})=G^{(N)}$: the target map commutes with the common projector,
which acts as the identity on $K$. More concretely, invariance gives
$zG^{(N)}=0$ and $eG^{(N)}=G^{(N)}$, while stable coefficients give
$\rho G^{(N)}=G^{(n)}$. The retained-vector identification and the argument
in Proposition~\ref{tx:lift-prop:lift} therefore yield $e\iota(G^{(n)})=G^{(N)}$. Hence invariant targets, including
$1$, are exactly preserved. For indexed targets this equality need
not hold. The image-subquotient argument in
Chapter~\ref{tx:chap:assemble-transfer} proves transfer of
prescribed indexed targets without assuming it.

\subsection{Axioms and joint tensor multiplication}
For an axiom index $a$ supported in $[n]$, let $F_n$ and $F_N$ be the
usable-axiom coefficient maps. Write $\iota_A$ for inclusion of retained
axiom indices and $e_A$ for the action of the same group-algebra element
on their source space. Naturality gives
\begin{equation}\label{tx:lift-eq:axiomlift}
             \Lift(f_{a,n})=F_N(e_A\iota_A(a)).
\end{equation}
The right side is a linear combination of actual usable axioms of
$\Gamma_N$. This is the precise sense in which axioms lift.

For the order example, a retained predecessor row has axiom
$\sum_{j<n}F_{ij}-1$. Plain inclusion omits the new summands with
$j$ in the added block, so it is not the larger row axiom.
Equation~\eqref{tx:lift-eq:axiomlift} repairs precisely this mismatch:
the lifted smaller row is a linear combination of the actual larger
axioms. The order, Boolean, functionality and predecessor-guard
equations are handled by the same usable-axiom map. We need no
injectivity of the predecessor function. The common projector fixes
the last line $1$, and joint tensor compatibility handles every
multiplication step, including one made possible by cancellation.

Now suppose $f\in V_j(n)\cap H_n$, so the smaller proof may derive
$x_vf$. Apply the same $e$ to the \emph{joint} tensor input and expand
by its variable coordinate:
\begin{equation}\label{tx:eq:joint-lift}
 e_{X_N\otimes H_N}\bigl(\iota(x_v)\otimes\iota(f)\bigr)
                  =\sum_{y\in\mathcal V_N}x_y\otimes g_y.
\end{equation}
Exactness and the natural tensor identification show that this vector
lies in $X_N\otimes(V_j(N)\cap H_N)$. Thus each $g_y$ belongs to
$V_j(N)$ and has degree at most $d-1$. Multiplication commutes with
the common group-algebra action, so
\begin{equation}\label{tx:lift-eq:liftmul}
                         \Lift(x_vf)=\sum_y x_yg_y.
\end{equation}
Each term is a permitted multiplication of a derivable polynomial.

There is even an explicit expression for these coefficients:
\[
            g_y=\sum_{g\in S_N:\,g(v)=y}a_g\,g\bigl(\iota(f)\bigr).
\]
Individual summands need not be known consequences; their grouped sum
$g_y$ is a consequence by the tensor argument. One must not replace the
joint operator in \eqref{tx:eq:joint-lift} by $e_X\otimes e_H$. That would
discard the correlations supplied by the common permutation $g$.

To produce a literal proof sequence, compute bases of the closure
stages while storing derivations for those bases. Express each $g_y$ in
the appropriate basis by linear algebra, append the multiplication
steps in \eqref{tx:lift-eq:liftmul}, and add. This is a finite procedure, with
no efficiency or proof-length claim. It explains exactly how
consequence-space lifting supplies a legal PC derivation.

\subsection{A multiplication step with all its proof lines}
The distinction between a joint action and two independent actions is
visible in a very small calculation. Work over any field, with three
labels $0,1,2$, variables $x_0,x_1,x_2$, and the uniform axioms
$\Gamma_3=\{x_0,x_1,x_2\}$. These mean $x_i=0$. Set
$d=2$, $X=\Span\{x_0,x_1,x_2\}$ and $H=P_{3,\le1}$.
Take $f=x_0+x_1$, let $\tau=(0\ 2)$ rename the labels, and put
$T=1-\tau$ in the group algebra. Applying $T$ to the joint input gives
\begin{align*}
 T_{X\otimes H}(x_0\otimes f)
 &=x_0\otimes(x_0+x_1)-x_2\otimes(x_2+x_1)\\
 &=x_0\otimes g_0+x_2\otimes g_2,
 \qquad g_0=x_0+x_1,\quad g_2=-x_2-x_1.
\end{align*}
Both grouped coefficients are derivable linear combinations of the
axioms. Multiplying the two tensor coordinates produces a literal
degree-two PC derivation:
\begin{center}
\begin{tabular}{@{}cll@{}}
\toprule
Line & Polynomial & Reason\\
\midrule
1 & $x_0$ & axiom\\
2 & $x_1$ & axiom\\
3 & $x_2$ & axiom\\
4 & $x_0+x_1$ & lines 1 and 2\\
5 & $-x_2-x_1$ & lines 3 and 2\\
6 & $x_0^2+x_0x_1$ & multiply line 4 by $x_0$\\
7 & $-x_2^2-x_2x_1$ & multiply line 5 by $x_2$\\
8 & $x_0^2+x_0x_1-x_2^2-x_2x_1$ & add lines 6 and 7\\
\bottomrule
\end{tabular}
\end{center}
Lines 4 and 5 have degree one, so their multiplication is legal within
the degree-two budget. The last line is $T(x_0f)$. There has been no
reduction modulo Boolean equations and no multiplication by a
polynomial of degree greater than one.

By contrast, acting independently on the two factors gives
\[
 (T_X\otimes T_H)(x_0\otimes f)
       =(x_0-x_2)\otimes(x_0-x_2),
\]
whose product is $(x_0-x_2)^2$. This is a different polynomial, also in
characteristic two: the joint result still has the terms involving
$x_1$. The common permutation in each summand is therefore essential.

This calculation illustrates the inference mechanism at one universe size.
The element $T$ is not a deletion projector and does not fix $1$.
For the actual larger-universe lift, replace $T$ by the common $e$
and use \eqref{tx:eq:joint-lift}. The grouped $g_y$ are then certified
consequences by exactness and the tensor identification, rather than
by the two-line calculation above. Once that certificate is available,
the final multiplication and addition steps are exactly the same.


\section{A one-label lifting example}\label{tx:lift-sec:small}
Work in characteristic two. Retain the label $a$ and add $b,c$.
Let $z=1+(bc)$ and $w=(ab)$ act on the linear variables
$x_a,x_b,x_c$. On their three-dimensional span,
\[
 z=\begin{pmatrix}0&0&0\\0&1&1\\0&1&1\end{pmatrix},\qquad
 w=\begin{pmatrix}0&1&0\\1&0&0\\0&0&1\end{pmatrix}.
\]
Direct multiplication gives $zwz=z$ and
\[
                       e=(1-zw)(1-wz)
                         =\begin{pmatrix}1&1&1\\1&1&1\\1&1&1\end{pmatrix}.
\]
Consequently
\begin{equation}\label{tx:lift-eq:smalllift}
                  \Lift(\alpha+\beta x_a)
                      =\alpha+\beta(x_a+x_b+x_c).
\end{equation}
Its retained coordinate is the original polynomial, and $\Lift(1)=1$.

Take the uniform degree-one axiom $\sum_{i<\ell}x_i=0$ at every size $\ell$.
The all-zero assignment satisfies this system at every size.
At size one the axiom is $x_a$; at size three it is $x_a+x_b+x_c$.
Formula~\eqref{tx:lift-eq:smalllift} sends the former exactly to the latter and
lifts every degree-one consequence. Plain inclusion would leave $x_a$,
which is not a consequence of the larger system: the assignment
$(x_a,x_b,x_c)=(1,1,0)$ satisfies its sum but not $x_a=0$.

Here the invariant target $G_\ell=\sum_{i<\ell}x_i$ is transported to
the prescribed larger target. In contrast, the indexed expression
$G_a^{(\ell)}=x_a$ has projector lift $x_a+x_b+x_c$, not $x_a$. The map
is linear but is not a ring homomorphism and does not preserve every
indexed expression literally. This calculation illustrates the distinction but lies outside the
indexed transfer range: for $k=d=t=1$ and $s=0$, where $t=1$ bounds
the support of the target index, that theorem uses
$R(d)=3$, whereas here $q=2$ and $n=1$. It is therefore no
counterexample to indexed periodicity.

The failure of a literal indexed lift requirement also has a general
explanation. At an admissible universe size, naturality for the invariant target
$\sum_{i<n}x_i$ gives
$\Lift(\sum_{i<n}x_i)=\sum_{i<N}x_i$. If every retained indexed
target $x_i$ had lift exactly $x_i$, linearity would instead give
$\sum_{i<n}x_i$, a different polynomial. Thus at least one indexed
target must change under the projector lift. This is why the proof
uses actual target-map images and nonzero-subquotient detection.
For the displayed invariant example,
$r=1$ and $q=2>r$, so a two-label block suffices. Larger
support budgets require larger blocks. The sufficient witness threshold
$n\ge2r-1$ is attained in this example; the general template theorem
uses a slightly more conservative onset.

\section{Descent to the coefficient field}\label{tx:lift-sec:descent}
Here \emph{descent} means returning a consequence established with
auxiliary coefficients to the original coefficient field.
The auxiliary choices have to meet the conditions of the block and
common-witness constructions; they are not probabilistic assumptions.
Their existence is supplied by those constructions.

The generic block coefficients may require an extension field. This
does not change the derivation degree of any target with coefficients in $K$: at fixed degree, scalar extension
commutes with the finite-dimensional spans, intersections and images
in \eqref{tx:eq:closure}. In particular,
\[
                       V_j(n;E)=E\otimes_K V_j(n;K),
\]
with the analogous equality for $C_d$. Membership of every $K$-coefficient target is therefore unchanged,
so each target-membership equivalence proved over $E$ holds over $K$.

There is also a concrete $K$-linear lifting map satisfying the forward
requirements \eqref{tx:eq:newcontract}. All coefficients used in the
construction lie in a finite extension $E/K$ inside the algebraic
closure. Choose a $K$-linear functional $\lambda:E\to K$ with
$\lambda(1)=1$, for example the coefficient of $1$ in a vector-space
basis containing $1$. Apply it coefficientwise to polynomials and put
\begin{equation}\label{tx:lift-eq:fieldlift}
                 \widehat\Lift(f)=\lambda_*\bigl(\Lift_E(f)\bigr)
                      \qquad(f\in W_n(K)).
\end{equation}
The finite-extension assertion uses only finitely many coefficients
of the block operator and group-algebra witness. Each is algebraic
over $K$ when chosen in its algebraic closure; adjoining this finite
list gives a finite-dimensional extension.

The functional $\lambda$ is a linear projection, not a field map.
For example, over $\mathbb F_2\subseteq\mathbb F_4$ with
$\omega^2+\omega+1=0$, take $\lambda(a+b\omega)=a$.
Then $\lambda(\omega)=0$, but $\lambda(\omega^2)=1$.
This explains why descent works through linear consequence spaces:
it is not substitution by a multiplicative map of coefficients.

This map preserves degree, satisfies $\rho\widehat\Lift=\mathrm{id}$
and fixes $1$. If $f\in V_j(n;K)$, then $\Lift_E(f)$ is an
$E$-linear combination of a $K$-basis of $V_j(N;K)$. Applying
$\lambda_*$ leaves a $K$-linear combination of that same basis.
Thus $\widehat\Lift(f)\in V_j(N;K)$, as required.

This is descent of consequence-space membership, not coefficientwise
projection of an arbitrary proof with extension-field scalars. Such a
projection need not preserve each scalar inference. A proof over $K$
is recovered from the $K$-closure bases. We use only the forward lifting
properties for $\widehat\Lift$; the stronger reflection of individual
consequences in Proposition~\ref{tx:lift-prop:lift} refers to the projector lift
over $E$. The two-way prescribed-target statement over $K$ follows separately
from target-map transport, nonzero-subquotient detection and field
invariance, as above.


\section{A summary of the map's three different uses}
For any smaller consequence $f$, the projector lift gives a larger
consequence $\mathsf L(f)$ of no greater degree. For a full stable
invariant target fitting the budget, this vector is exactly the
prescribed larger target $G^{(N)}$. For general indexed targets, literal
pointwise equality is not available: the actual target-map image and
the quotient $(C_d(N)+\im J_N)/C_d(N)$ supply the transfer instead.
For NS, use the analogous quotient by $I_d(N)$. These quotients test
whether a target lies outside the relevant consequence space.

These uses share one construction but have different logical inputs.
Keeping them separate is the simplest way to avoid proving only a
statement about an unspecified lifted vector while claiming a theorem
about the specified target.

\begin{checkpoint}
A single common projector makes the entire closure calculation compatible.
It does not replace ordinary-PC multiplication by a new proof rule.
The larger consequence remains derivable by the original legal rules.
\end{checkpoint}
